\section*{Capítulo \ref{ch:g3:addsub} --- Suma y resta}

\begin{solution}{exo:g3:addsub:1}
$346 + 253 = 599$; \quad
$478 + 265 = 743$; \quad
$1\,536 + 2\,876 = 4\,412$.
\end{solution}

\begin{solution}{exo:g3:addsub:2}
$587 - 234 = 353$ (controlar: $353 + 234 = 587$).

$805 - 348 = 457$ (controlar: $457 + 348 = 805$).

$1\,000 - 463 = 537$ (controlar: $537 + 463 = 1\,000$).
\end{solution}

\begin{solution}{exo:g3:addsub:3}
$38 + 7 = 38 + 2 + 5 = 45$; \quad
$56 + 9 = 56 + 4 + 5 = 65$; \quad
$45 + 28 = 45 + 30 - 2 = 73$.
\end{solution}

\begin{solution}{exo:g3:addsub:4}
$340 + 220 = 560$; \quad $675 + 210 = 885$; \quad
$512 - 300 = 212$; \quad $840 - 220 = 620$.
\end{solution}

\begin{solution}{exo:g3:addsub:5}
De $47$ a $50$: $3$; de $50$ a $63$: $13$; total
$3 + 13 = 16$. Las columnas confirman: $63 - 47 = 16$.
\end{solution}

\begin{solution}{exo:g3:addsub:6}
$492 + 218$: estimar $500 + 200 = 700$; exactamente $710$.

$913 - 388$: estimar $900 - 400 = 500$; exactamente $525$.
\end{solution}

\begin{solution}{exo:g3:addsub:7}
$348 + 296 = 644$. la escuela tiene $644$ estudiantes.
\end{solution}

\begin{solution}{exo:g3:addsub:8}
$224 - 178 = 46$. Nora todavía tiene $46$ páginas para leer.
\end{solution}

\begin{solution}{exo:g3:addsub:9}
Unidades: $7 + \text{?} = 2$ es imposible, entonces $7 + 5 = 12$: las unidades
el agujero es $5$, llevar $1$. Decenas: $\text{?} + 2 + 1 = 6$: el agujero de las decenas
es $3$ (sin transporte). Cientos: $3 + \text{?} = 8$: el agujero de las centenas es
$5$. Entonces $337 + 525 = 862$, y cada hoyo tenía exactamente un posible
dígito --- solo hay una manera.
\end{solution}

\begin{solution}{exo:g3:addsub:10}
Vendido: $145 + 89 = 234$ croissants. Horneado: vendido más no vendido,
$234 + 17 = 251$ croissants.
\end{solution}

\begin{solution}{exo:g3:addsub:11}
Tiene razón: $500 - 200 = 300$, más $1$ da $301$; en columnas,
$500 - 199 = 301$ también. quitando $200$ elimina \emph{uno demasiados},
entonces uno debe ser devuelto --- restando un vecino amigo y
corregir es a menudo más rápido que pedir prestado.
\end{solution}
